% ---------- Слайд 16: Заключение ----------
\begin{frame}[fragile,label=conclusion]{Заключение: данный доклад подготовлен в \LaTeX}
  \begin{columns}[T,onlytextwidth]
    % --- ЛЕВАЯ КОЛОНКА: КОД ---
    \column{0.52\textwidth}
      \textbf{Исходный код:}
      \vspace{0.2em}
      \inputminted[fontsize=\scriptsize, frame=single, breaklines=true]{latex}{examples/conclusion.tex}

    % --- ПРАВАЯ КОЛОНКА: РЕЗУЛЬТАТ ---
    \column{0.44\textwidth}
      \textbf{Результат:}
      \vspace{0.2em}
      \begin{block}{Заключение}
  Данный доклад полностью подготовлен
  в \LaTeX{} с использованием класса beamer.
\end{block}
  \end{columns}

  \vspace{1em}
  \centering
  \textbf{Исходные файлы презентации доступны по QR-коду:}\\
  \vspace{0.3em}
  % Замените ссылку на ваш реальный репозиторий
  \qrcode[height=1.3cm]{https://github.com/username/latex-presentation-template}
  
  \note{Наглядно продемонстрировать, что весь дизайн, формулы,
  анимации и таблицы на слайдах были созданы с помощью
  исходного кода.}
\end{frame}

% ---------- Слайд 17: Вопросы ----------
\begin{frame}[label=questions]{Вопросы}
  \begin{center}
    \Large
    \textbf{Благодарю за внимание!}\\
    \vspace{1em}
    \normalsize
    \textbf{Докладчик:} Александр Лихач\\
    \textbf{Email:} smouzi45@gmail.com
  \end{center}
  \note{Завершить доклад и перейти к вопросам аудитории.}
\end{frame}